% Place this subsection at the end of the Experiments section, immediately
% before Conclusion. Replace every \placeholder after collecting logs.
\subsection{Onboard Deployment Feasibility}
\label{sec:onboard_deployment}

We further integrated DAC-SFC into an onboard receding-horizon planning stack
to assess deployment feasibility.  The original EGO-Planner-v2 map, finite
local target, trajectory server, and flight-control interfaces were retained.
For each replanning request, a collision-free A* path was shortened by
farthest-visible connections and divided into segments of at most
$1.0\,\mathrm{m}$.  The resulting polyline was passed to the same unoptimized
MINCO deformation probe, CSGN construction, Active-Witness corridor generator,
and GCOPTER backend used in the main experiments.  The generated quintic
trajectory was converted directly to the existing trajectory-server format.
An independent check against the inflated local occupancy map was required
before trajectory activation; planning failures retained the original EGO
planner as a fallback.

The purpose of this study is to verify online execution and interface
compatibility rather than to provide another algorithm comparison.  We record
the A* search, route-shortening, deformation-probe, CSGN, corridor-generation,
optimizer-setup, and optimization times, together with the number of corridor
facets and geometric face evaluations per objective/gradient call.  During
development away from the vehicle, the existing simulator supplied the depth,
pose, and odometry interfaces.  The same logging and safety gates were then
used unchanged for onboard trials.

Across \placeholder{N} onboard replanning attempts, the complete DAC-SFC
pipeline succeeded in \placeholder{success-rate}\% of the attempts, with
\placeholder{fallback-count} fallback activations.  Median and 95th-percentile
end-to-end planning latency were \placeholder{median-total} and
\placeholder{p95-total}\,ms, respectively.  The median optimization latency was
\placeholder{median-opt}\,ms, while the generated corridors contained a median
of \placeholder{median-faces} facets, corresponding to
\placeholder{median-geo-evals} geometric face evaluations per
objective/gradient call.  All activated trajectories passed the independent
inflated-map check; the observed maximum velocity and acceleration were
\placeholder{max-velocity}\,$\mathrm{m\,s^{-1}}$ and
\placeholder{max-acceleration}\,$\mathrm{m\,s^{-2}}$.  These results demonstrate
that DAC-SFC can be executed within the onboard replanning loop using the
existing sensing and control interfaces.
